\section{Back to the constant}\label{end:visible}

The first construction fitted a list of arithmetic questions into one
real number. Each prime carried a question; its lane repeated the answer.
That repetition let a single missing answer alter a limiting percentage,
break a self-similarity, and turn Euler's convergent sum into infinity.

Reading the positions through coprimality changed the questions we could
ask. Correlations counted missing lanes. Their generating functions
turned that count into degrees, zeros and spectral directions. A random
scale connected two ways of averaging, and its density became the curve
whose derivatives recognize primes. The prime and composite parts could
even be put back together to recover an ordinary exponential law.

Then we hid the numbers. The relationships still revealed prime names.
The cancellation of three observed frequencies explained why; the noisy
experiment quantified how far the idea could go. Its paradox survives
the precise rules of the game: the exact shared primes of a fraction of
pairs tending to one can be recovered while the retained fraction of
information about the complete table tends to zero. The surviving
information grows, but the uncertainty about the whole table grows
faster. Both facts belong in the picture.

The later experiments asked what arithmetic arrangements look like.
Coprime crowds could leave earlier holes in the list of covered primes.
Averaging a matrix and averaging its characteristic polynomial revealed
different spectra. Conditioning a factorization to be unusually balanced
produced a Gaussian profile. These conclusions make sense on their own;
the shared constructions gave us a route from one question to the next.

Finally we returned to the original assignment and removed seven.
Almost all primes survived every fixed number of rounds. Only three
survived forever. The same arithmetic could appear abundant or almost
empty, depending on when we looked.

A last change of view sent the curve into the complex plane. Its prime
coefficients made it meet every polynomial infinitely often, yet prevented
any three points with the same value from forming an equilateral triangle
centred at the origin. A rotation that keeps just one prime explained both
conclusions.

Section~\ref{open:section} gathers the conjecture and questions left
by these experiments: their hidden messages, complex values, spectra
and long-term behavior.  The earlier connections to classical
arithmetic questions also remain invitations to further work.
Goldbach, RH and twin primes remain open, and a new formulation
needs new input before it can settle them.

The unifying question is simpler: what can a number remember after we
change the way we observe it? Here it remembered arithmetic in a digit
frequency, in a curve's missing derivatives, in the signs of a matrix,
and in noisy relationships between closed boxes. Each observation
revealed something the previous one had made difficult to see.
